\documentclass[aps,pre,onecolumn,10pt,superscriptaddress]{revtex4-2}
\usepackage{amsmath,amssymb,amsthm}
\usepackage[hidelinks]{hyperref}
\newcommand{\norm}[1]{\left\lVert #1\right\rVert}
\newcommand{\RMS}{\operatorname{RMS}}
\begin{document}
\title{Supporting calculations: loss accuracy and balanced norm shape}
\author{Govind Gandhi}
\affiliation{School of Data Science, University of Virginia, Charlottesville, Virginia 22904, USA}
\author{Sadamori Kojaku}
\affiliation{School of Systems Science and Industrial Engineering, Binghamton University, Binghamton, New York 13902, USA}
\author{Yong-Yeol Ahn}
\affiliation{School of Data Science, University of Virginia, Charlottesville, Virginia 22904, USA}
\date{\today}
\maketitle

These calculations concern a fixed ordered SGNS objective with separate center and context factors.
For \(X=UV^\top\in\mathbb R^{n\times n}\), write
\begin{equation}
 L(X)=\sum_{ij}\{a_{ij}\log(1+e^{-X_{ij}})+b_{ij}\log(1+e^{X_{ij}})\},
 \qquad r_i(X)^2=[(XX^\top)^{1/2}]_{ii}.
\end{equation}
The canonical balanced center radii \(r_i\) depend only on the score matrix.
Set \(\rho(X)=\RMS(r)=\sqrt{n^{-1}\sum_i r_i^2}\) and \(\widetilde r=r/\rho\) when \(\rho>0\).
Balancing removes factorization gauge; RMS normalization removes overall radius scale.
Neither operation implies that near-optimal loss identifies a unique normalized profile.

\section{Exact sparse-graph construction}

Take a triangle on nodes \(0,1,2\), with a pendant node \(3\) attached to \(0\).
Its degree vector is \(d=(3,2,2,1)\) and its volume is eight.
Fix uniform center exposure, one-step positive contexts, stationary noise \(q_j=d_j/8\), and five collision-allowing negatives:
\begin{equation}
 a_{ij}=\frac{A_{ij}}{4d_i},\qquad b_{ij}=\frac{5d_j}{32}.
\end{equation}
On an edge, the scalar loss is minimized at the symmetric score
\(X_{ij}^\star=\log[8/(5d_id_j)]\).
On a zero-positive pair, including each diagonal, its infimum is zero as the score tends to \(-\infty\).
The total negative weight of these pairs is \(B_0=35/16\).

Given any tolerance \(\delta>0\), put \(s=\max\{0,\log(B_0/\delta)\}\).
Keep every edge score at its optimum and set every off-diagonal nonedge score to \(-s\).
Let \(R_i\) denote the absolute off-diagonal row sum of this common symmetric matrix.
For any finite collection of strictly positive desired node profiles \(h^{(p)}\), set
\begin{equation}
 \alpha_i^{(p)}=\frac{(h_i^{(p)})^2}{\RMS(h^{(p)})^2},
 \qquad
 D=\max_{p,i}\frac{\max\{s,R_i+1\}}{\alpha_i^{(p)}},
 \qquad X_{ii}^{(p)}=-D\alpha_i^{(p)}.
\end{equation}
Each \(X^{(p)}\) is strictly negative definite by symmetric diagonal dominance.
Consequently \((X^{(p)}X^{(p)\top})^{1/2}=-X^{(p)}\), and hence
\begin{equation}
 r_i(X^{(p)})^2=D\alpha_i^{(p)},\qquad
 \rho(X^{(p)})=\sqrt D,\qquad
 \widetilde r_i(X^{(p)})=\frac{h_i^{(p)}}{\RMS(h^{(p)})}.
\end{equation}
All positive-pair terms attain their scalar minima, while every zero-positive score is at most \(-s\).
Thus, for the common unrestricted infimum \(L_{\inf}\),
\begin{equation}
 0<L(X^{(p)})-L_{\inf}
 \le B_0\log(1+e^{-s})\le B_0e^{-s}\le\delta.
\end{equation}
This includes \(\delta\ge B_0\), when \(s=0\).
Explicit balanced factors are \(U=(-X)^{1/2}\) and \(V=-(-X)^{1/2}\); four dimensions suffice.
Choosing degree-level profiles \((1,2,3)\), \((3,2,1)\), and \((1,3,1)\) at degrees \(1,2,3\) gives increasing, decreasing, and peaked curves.
The RMS uses all four nodes, including both degree-two nodes.
The profiles share one RMS at each tolerance; that RMS can grow as \(\delta\) decreases.
Their losses need not be exactly equal, and the construction does not claim that an optimizer selects all three shapes.
It does not establish nonidentification for the rank-24 fits on the larger experimental graphs.

\section{A conservative full-support bound}

Suppose now that every \(a_{ij},b_{ij}>0\).
The unique unrestricted optimum is \(X^\star_{ij}=\log(a_{ij}/b_{ij})\).
Define
\begin{equation}
 \Delta=L(X)-L(X^\star),\quad
 h=\min_{ij}\frac{a_{ij}b_{ij}}{a_{ij}+b_{ij}}>0,\quad
 \psi(t)=t-1+e^{-t},\quad E=\norm{X-X^\star}_{\mathrm F}.
\end{equation}
For the scalar loss \(\ell(x)=a\log(1+e^{-x})+b\log(1+e^x)\),
\(\ell''(x)=(a+b)e^x/(1+e^x)^2\) and \(\lvert(\log\ell'')'(x)\rvert\le1\).
Integrating this curvature lower bound away from \(x^\star\) gives
\(\ell(x^\star+z)-\ell(x^\star)\ge\ell''(x^\star)\psi(|z|)\).
Moreover,
\begin{equation}
 \frac{\psi(t)}{t^2}=\int_0^1(1-u)e^{-ut}\,du
\end{equation}
is decreasing for \(t\ge0\), with the continuous extension at zero.
Since each \(|z_{ij}|\le E\), summing yields \(\Delta\ge h\psi(E)\).
The inequality \(\psi(t)\ge t^2/(t+2)\) follows because
\((t+2)\psi(t)-t^2=t-2+(t+2)e^{-t}\) vanishes at zero and has nonnegative derivative \(1-(t+1)e^{-t}\).
Therefore
\begin{equation}
 E\le\psi^{-1}(\Delta/h)\le B,
 \qquad B=\frac{c+\sqrt{c^2+8c}}2,\qquad c=\Delta/h.
 \label{eq:score-bound}
\end{equation}

To pass from scores to radii, form the Hermitian dilation
\(H(X)=\left(\begin{smallmatrix}0&X\\X^\top&0\end{smallmatrix}\right)\).
Its absolute value has diagonal blocks \((XX^\top)^{1/2}\) and \((X^\top X)^{1/2}\).
For Hermitian matrices \(A,B\), expansion in their two eigenbases and
\((|\lambda|-|\mu|)^2\le(\lambda-\mu)^2\) imply
\(\norm{|A|-|B|}_{\mathrm F}\le\norm{A-B}_{\mathrm F}\).
Applying this to the dilations gives
\begin{equation}
 \norm{(XX^\top)^{1/2}-(X^\star X^{\star\top})^{1/2}}_{\mathrm F}
 \le\sqrt2 E.
\end{equation}
Using \((r_i-r_i^\star)^2\le|r_i^2-r_i^{\star2}|\) and Cauchy--Schwarz,
\begin{equation}
 \RMS(r-r^\star)\le e,\qquad e=(2/n)^{1/4}\sqrt B.
\end{equation}
For nonzero fitted and target RMS, the reverse triangle inequality gives
\begin{equation}
 \RMS(\widetilde r-\widetilde r^\star)
 \le\min\{\sqrt2,\,2e/\rho^\star\}.
\end{equation}
The cap \(\sqrt2\) uses nonnegative unit-RMS profiles; it is a trivial bound, not a useful loss-based certificate.
If \(e<\rho^\star\), nonzero fitted RMS is guaranteed.
A zero target RMS makes its normalized shape undefined.
For a bin of \(m\) nodes, its normalized mean error is at most \(\sqrt{n/m}\) times the profile RMS bound.

The bound concerns distance to the unrestricted optimum, including for a rank-constrained fitted matrix.
It neither proves uniqueness of a rank-constrained optimum nor separates approximation error from incomplete optimization when \(\Delta\) is large.
Small positive pair weights can make \(h\) tiny and the certificate uninformative.
Failure to certify a curve is not evidence that it is unstable.
The accompanying condition table reports this limitation alongside actual shape differences; it does not convert repeated-start agreement into a uniqueness claim.

\end{document}
